% Part: computability
% Chapter: computability-theory
% Section: universal-part-function

\documentclass[../../../include/open-logic-section]{subfiles}

\begin{document}

\olfileid[hi]{cmp}{thy}{uni}
\olsection{सार्वत्रिक आंशिक संगणनीय फलन}

\begin{thm}
\ollabel{thm:univ-comp}
एक सार्वत्रिक आंशिक संगणनीय फलन~$\fn{Un}(e,x)$ है। दूसरे शब्दों में,
एक फलन $\fn{Un}(e,x)$ ऐसा है कि:
\begin{enumerate}
\item $\fn{Un}(e,x)$ आंशिक संगणनीय है।
\item यदि $f(x)$ कोई आंशिक संगणनीय फलन है, तो कोई प्राकृतिक संख्या
$e$ ऐसी है कि $f(x) \simeq \fn{Un}(e,x)$ प्रत्येक~$x$ के लिए।
\end{enumerate}
\end{thm}

\begin{proof}
$\fn{Un}(e,x) \simeq U(\umin{s}{T(e,x,s)})$ मानें, जहाँ $U$ और $T$,
क्लीन के प्रसामान्य रूप प्रमेय (\olref[nfm]{thm:normal-form}) के अनुसार हैं।
\end{proof}

\begin{explain}
यह केवल इस बात को सटीक रूप से कहने का ढंग है कि हमारे पास आंशिक संगणनीय
फलनों का एक प्रभावी परिगणन है। आशय यह है कि यदि $f_e$ से उस फलन को
लिखें जिसकी परिभाषा $f_e(x) = \fn{Un}(e,x)$ है, तो अनुक्रम $f_0$,
$f_1$, $f_2$, \dots में सभी आंशिक संगणनीय फलन आते हैं और $f_e(x)$ को
$e$ तथा~$x$ में ``एकसमान रूप से'' संगणित किया जा सकता है। सरलता के लिए
हम ऐसे द्विस्थानी फलन का उपयोग कर रहे हैं जो एकस्थानी फलनों के लिए
सार्वत्रिक है, पर संख्याओं के अनुक्रमों को कूटित करके इसे अधिक तर्कों के
लिए सहज ही सामान्यीकृत किया जा सकता है। उदाहरणतः, ध्यान दें कि यदि
$f(x,y,z)$ कोई $3$-स्थानी आंशिक पुनरावर्ती फलन है, तो फलन
$g(x) \simeq f((x)_0, (x)_1, (x)_2)$ एकस्थानी पुनरावर्ती फलन है।
\end{explain}

\end{document}
